\section{Conclusion}

This paper examined which components limit STAMP in open-world test-time
adaptation under severe image corruption. We tested four focused interventions
while retaining the consistency filter, class-balanced replay memory, and
shared CIFAR-C evaluation protocol. Replacing entropy with KNN feature distance
improved known-class accuracy but did not improve the joint
recognition--rejection trade-off. Energy scoring substantially improved the
CIFAR-100-C result, but its CIFAR-10-C AUC remained weak. Applying ACI to the
energy threshold produced nearly identical, slightly lower metrics than the
rolling-percentile energy baseline on both datasets; hence, the reported
experiment does not support adaptive thresholding as the main remedy in these
streams.

The most consequential intervention was correcting replay-loss scaling while
memory is partially filled. STAMP+BSN achieved the best H-scores: 81.19 on
CIFAR-10-C with noise outliers and 72.89 on CIFAR-100-C.
These results are consistent with the hypothesis that excessive early replay
updates caused by the $64/N$ scaling error can destabilize adaptation, although
the aggregate metrics alone do not prove this mechanism. Overall, the findings
shift the emphasis from selecting an outlier score alone to ensuring that the
memory-adaptation objective is correctly normalized.

The study is limited to severity-5 CIFAR-C corruption streams with generated
noise outliers and reports aggregate means without multi-seed uncertainty,
per-corruption metrics, or memory-admission traces. Future work should log
replay occupancy, loss scale, ID rejection, and OOD admission over time, and
test the replay correction across multiple seeds, semantic OOD sources, and
larger open-world benchmarks.
